\section{讨论、局限与研究工程经历}
\subsection{反馈必须对应到实际评价目标}
两个案例都说明，过程中的反馈量与最终目标需要分别解释。PPO 的训练奖励来自不同题目和随机采样，任务变化应由同条件配对验证衡量；程序搜索中的漂亮解释或紧密图形，也必须经过原评估器执行后才算有效候选。本文保留生成失败、无效候选和缺失结果，避免只呈现成功记录。

参数优化案例回答了“能否走通并验证”的问题：完整训练、保存、Actor 重载和配对评估均已完成。程序搜索案例则展示了有限记录中的有效生成与选择，但全局最佳来自初始化，不能把“进行了演化”直接写成“每一代都带来提升”。两个任务的分数不可互相换算，也没有构成方法间公平竞赛。

\subsection{哪些限制仍然存在}
\begin{itemize}
\item \key{实验规模。}参数优化只使用一个种子、一个短算术任务和一个开发验证划分，没有多种子、独立测试集或消融比较。净增加五道正确题是描述性结果，不证明统计显著、最佳配置或通用推理能力提升。
\item \key{协议形成过程。}早期输出参与过首行答案规则的修订，因此不把当前验证集描述为从未查看的测试集。规则确定后的新 A10 前后比较保持一致，但不能消除整个开发过程的选择影响。
\item \key{保存范围。}Actor 实际重载已验证；Critic 只有适配器权重检查。其导出包含 LoRA 矩阵，却缺少底层随机初始化的价值头权重，因此仅凭这些文件无法重建完整的已训练价值函数。完整优化器恢复、长期续训和独立 Critic 重建未得到验证。
\item \key{搜索预算与终止。}圆装填仅五次有效评估，使用者在导出后手动停止。第 2 代开始不等于预算自动完成，也不能代表官方更大规模基准。
\item \key{可复现范围。}导出的代码和数据可复核，外部模型服务的生成响应不保证逐次相同。在线与本地数值差异的原因也未被单因素实验隔离。
\end{itemize}

\subsection{从失败恢复中形成的工程认识}
早期 T4 运行完成了 32 批计算，却在最终保存时因 PEFT 权重别名检查失败。这说明“训练循环结束”“模型成功保存”“重载后仍可使用”是不同证据。后续先做独立的 16 题保存重载检查，再进行完整运行，能够更早发现交付环节的问题；历史失败证据仍保留在附录~\ref{app:history-zh}，没有用新分数覆盖旧记录。

11 月 8 日的云端网络问题还阻碍了 FlashAttention 下载。源码检查确认当前单 GPU、未打包的 PPO 路径不会调用相关环形注意力功能，于是用哈希受约束的补丁延后三处导入。真正调用相关功能时仍需要真实依赖，没有实现假的算子。原依赖声明未篡改，\code{pip check} 中这一项缺失也被明确记录。此前的实际 CUDA Actor/Critic 前向反向检查和 FusedAdam 参数更新，随后完整运行，共同支持受限路径可用。

PPO 明确使用 eager 注意力；前后评估保留同一份原函数，可能使用 SDPA，没有强制改变实现。这里的工作是兼容性修复与运行验证，不是提出新的 PPO 算法，也没有把注意力实现选择当作性能改进的独立原因。

程序搜索早期有四次已完成的空代码回复，修正版提高输出上限并减少预算后才取得有效评估。思考耗尽输出额度是合理猜测，\cite{deepseek} 但同时改变了不止一个设置，也缺少完整结束原因，不能确认为单一根因。第 1 代日志报告的五次代码实现平均约 5 分 38 秒，评估平均约 30 秒；因此等待并不都来自几何计算。Planner 的 70,376 tokens 仅覆盖七个保存方案，不能冒充全任务成本。

\section{后续改进与可检验方案}
\subsection{以评估器证据约束反思和记忆}
重复的非法程序会浪费模型调用。可把重叠圆编号、越界、数组形状错误和超时整理为结构化反馈，附上相关代码，要求提出具体修复假设，再由原评估器验证。

平台已有错误反思和分层记忆。\cite{memory} 改进重点是更严格地筛选经验：每条经验关联候选、实测分数和评估器版本；区分已验证事实与猜测；移除重复、过期或矛盾内容；每次只检索少量相关记录。例如，经检查能够保持约束的半径更新值得保存，自称高分却越界的程序不能成为成功经验。

预期收益是减少重复错误、提高合法候选比例，但反思与检索会消耗 tokens，陈旧经验也可能把搜索带偏。验证时固定模型、评估器和总预算，比较不使用这类记忆与使用已验证记忆的运行；跨多个种子记录合法率、重复失败率、最佳有效分数和成本。本文没有执行这组对照实验。

\subsection{保留搜索多样性，并依据实际进展分配预算}
若每个提示都模仿当前最佳代码，搜索可能重复同一方向。可给不同岛分配不同提示角色：局部微调当前布局、探索新空间排列、优先生成可靠几何构造。各岛仍使用同一目标与评估器，完全重复的代码先去重，候选只定期交换。

分岛与迁移本来就是框架机制。\cite{islands} 进一步的研究问题是，能否依据有效分数、停滞和失败率分配后续调用：某方向长期无进展时，把部分预算留给替代方向。模型可以建议切换，但最终资源分配依据应是实际评估，而不是自然语言解释是否新颖。

多样性可能帮助离开局部最优，自适应分配也可能减少重复调用；代价是探索弱方向和短运行的不稳定。公平实验应固定总调用量或 token 成本，对比统一提示、固定分配与多角色提示、自适应分配；用多个种子报告最佳有效分数的中位数和分布、独特合法候选数及达到预设目标的耗时。这里仍是待验证方案，没有声称已经获得改进。

\smallheading{共同的评估边界}
评估器应位于模型不可修改的区域，候选程序最终独立复核，失败也应保留。对于单个 26 圆实例，结论只限该实例；改变问题规模或任务后仍需重新测试。两种提案都不必先重训大语言模型，却都需要真实对照结果才能成为有效方法的证据。

\section{结论}
本研究完成了低秩适配实现验证，以及Qwen 基座上的完整小规模 PPO 流程。在固定开发验证条件下，正确数由 117/128 变为 122/128；Actor 可保存并重载，Critic 适配器通过权重检查。第二个案例以真实在线调用生成并评估了五个合法圆装填候选，最佳在线分数为 2.62001380699489，独立本地复现记录另行保存。两项研究支持有限任务下的实现可行性与可检查结果，不提供普遍能力、全局最优或新算法的证明。

对后续研究而言，主要尚未回答的问题是不同种子、任务和预算下的稳定性，以及两项代理改进能否在相同成本下减少无效搜索。原始程序、回答、几何数据、失败记录和源码版本为这些进一步实验提供了起点。
